\documentclass[12pt]{article}
\newcommand{\bandname}{Grades 2--3}
\newcommand{\packetid}{FNN-23-v1}
\usepackage[letterpaper, left=1.4cm, right=1.4cm, top=2.0cm, bottom=2.0cm,
  headheight=16pt, headsep=0.45cm, footskip=0.85cm]{geometry}
\usepackage[T1]{fontenc}
\usepackage{tgheros}
\renewcommand{\familydefault}{\sfdefault}
\usepackage{tikz}
\usepackage{pifont}
\usepackage{fancyhdr}
\pagestyle{fancy}
\fancyhf{}
\lhead{\normalsize Week NN / Graceful labels / \bandname}
\lfoot{\normalsize Bellingham Math Circle / Week NN / \packetid}
\rfoot{\normalsize \thepage}
\renewcommand{\headrulewidth}{0pt}
\renewcommand{\footrulewidth}{0pt}
\setlength{\parindent}{0pt}
\setlength{\parskip}{0pt}
\newcommand{\problem}[1]{\textbf{Problem #1:}}
\newcommand{\fig}[1]{\input{figs/#1.tex}}
\raggedright

\begin{document}
\setlength{\baselineskip}{17pt}

Put a different number card on each dot, choosing from the cards named under the network. In the box on each link, write the difference between the numbers at its two ends. The network is labelled when its link numbers are 1, 2, 3 and so on up to the number of links, each once.

\medskip
\begin{center}
\fig{rules_example}
\end{center}

\vspace{0.5cm}
\problem{1} Label each path.

\vspace{0.4cm}
\begin{center}
\fig{path4}

\vspace{0.7cm}
\fig{path5}

\vspace{0.7cm}
\fig{path6}
\end{center}

\newpage
\problem{2} Label this path.

\vspace{0.8cm}
\begin{center}
\fig{path11}
\end{center}

\newpage
\problem{3} Label each network.

\vspace{0.8cm}
\begin{center}
\fig{tree_fork}

\vspace{2.0cm}
\fig{tree_h}
\end{center}

\newpage
\problem{4} Which cards can go in the middle dot when you label this network? Explain why the other cards cannot.

\vspace{0.6cm}
\begin{center}
\fig{star5}

\vspace{1.0cm}
\fig{star5_small}
\end{center}

\newpage
\problem{5} Label each loop, or cross it out if it cannot be done.

\vspace{0.6cm}
\fig{loop3}\hfill\fig{loop4}

\vspace{0.8cm}
\raisebox{0.685cm}{\fig{loop5}}\hfill\fig{loop6}

\newpage
\problem{6} Label this loop.

\vspace{0.6cm}
\begin{center}
\fig{loop7}
\end{center}

\vspace{1.0cm}
\problem{7} Explain how you know that the loops you crossed out in Problem 5 cannot be labelled.

\newpage
\problem{8} Put four counters on the strip, one on 0 and one on 6, so that the six distances between pairs of counters are 1, 2, 3, 4, 5 and 6.

\vspace{0.5cm}
\begin{center}
\fig{strip_example}

\vspace{1.6cm}
\fig{strip6_big}

\vspace{1.6cm}
\fig{strip6_small}

\vspace{0.9cm}
\fig{strip6_small}

\vspace{0.9cm}
\fig{strip6_small}
\end{center}

\end{document}
